\documentclass[a4paper,10pt]{article}

\usepackage[utf8]{inputenc}
\usepackage[T1]{fontenc}
\usepackage[margin=1.5cm]{geometry}
\usepackage{enumitem}
\usepackage{hyperref}
\usepackage{multicol}
\usepackage{titlesec}

\hypersetup{colorlinks=true, urlcolor=black, linkcolor=black}

\titleformat{\section}{\large\bfseries}{}{0em}{}[\titlerule]
\titlespacing{\section}{0pt}{6pt}{4pt}

\setlength{\parindent}{0pt}
\setlength{\parskip}{2pt}
\renewcommand{\familydefault}{\sfdefault}

\begin{document}

% =================================================================
% BLOCO 1: CABEÇALHO
% =================================================================
% OBJETIVO: Identificação imediata. O ATS extrai nome, contato e
% título diretamente daqui. Mantenha tudo em texto puro, sem tabelas
% ou caixas --- ATS não lê formatações complexas.
%
% REGRAS:
%   - Nome completo em destaque (sem abreviações)
%   - Título profissional: use as palavras-chave EXATAS da vaga.
%     Ex: se a vaga pede "Desenvolvedor Java Sênior", use isso,
%     não "Engenheiro de Software Backend"
%   - Telefone, e-mail, LinkedIn e GitHub em texto simples
%   - Não coloque foto, CPF, data de nascimento ou endereço completo
%   - LinkedIn e GitHub devem ter URLs clicáveis (hyperref faz isso)
%
% DICA ATS: O título (linha 2) é o campo mais lido pelos filtros
% automáticos. Espelhe o headline da vaga palavra por palavra.
% =================================================================

\begin{center}
{\LARGE \textbf{SEU NOME COMPLETO}} \\[3pt]
\textbf{Título Profissional $|$ Tecnologia 1 $|$ Tecnologia 2 $|$ Tecnologia 3 $|$ Cloud $|$ CI/CD} \\[4pt]
\small
País $\bullet$ Disponível para mudança $\bullet$ Remoto / Híbrido / Presencial \\
(XX) XXXXX-XXXX $\bullet$ seuemail@email.com \\
\href{https://www.linkedin.com/in/luisfernando-eng}{linkedin.com/in/luisfernando-eng}
$\bullet$
\href{https://github.com/seu-usuario}{github.com/seu-usuario}
\end{center}



% =================================================================
% BLOCO 2: RESUMO PROFISSIONAL
% =================================================================
% OBJETIVO: É o primeiro texto que o recrutador lê após o ATS
% aprovar o currículo. Deve responder em 4 linhas: quem você é,
% o que você faz, com o quê e qual seu diferencial.
%
% REGRAS:
%   - 3 a 5 linhas, sem bullet points
%   - Inclua as principais palavras-chave da vaga em negrito
%   - Não use frases vagas: "profissional dedicado", "perfil dinâmico"
%   - Mencione linguagem principal, stack, tipo de sistema e contexto
%     (ex: alta volumetria, produção, fintech, e-commerce)
%   - Termine com seu diferencial: o que te separa dos outros 100
%     candidatos com o mesmo stack?
%
% ESTRUTURA SUGERIDA:
%   [Formação] com [X anos] de experiência em [linguagem/stack],
%   atuando em [tipo de sistema]. Experiência com [padrões/arquitetura].
%   [Diferencial concreto: projeto real, escala, resultado mensurável].
%   Perfil [colaborativo/hands-on] com foco em [qualidade/performance].
%
% DICA ATS: Repita aqui as mesmas palavras-chave do título.
% ATS pondera frequência --- palavra que aparece 2x tem mais peso.
% =================================================================

\section*{Resumo Profissional}

[Sua formação] com experiência no desenvolvimento de [tipo de aplicação] em \textbf{[Linguagem principal]} com \textbf{[Framework]}, atuando na construção e evolução de \textbf{[tipo de arquitetura]}. Experiência com \textbf{[padrão arquitetural]} e \textbf{[mensageria/evento]}, aplicando padrões como [padrões relevantes da vaga]. [Diferencial concreto: projeto real em produção, escala, impacto mensurável]. Perfil colaborativo, com participação em code reviews e aplicação de [boas práticas relevantes].


% =================================================================
% BLOCO 3: COMPETÊNCIAS TÉCNICAS
% =================================================================
% OBJETIVO: Este é o bloco mais importante para o ATS.
% Sistemas ATS varrem este campo em busca de palavras-chave
% que batem com os requisitos da vaga.
%
% REGRAS:
%   - Liste EXATAMENTE as tecnologias pedidas na vaga
%   - Use os nomes oficiais: "Spring Boot" (não "Springboot"),
%     "PostgreSQL" (não "Postgres"), "GitHub Actions" (não "GH Actions")
%   - Organize por categoria (Linguagens, Backend, Cloud etc.)
%   - Não liste tecnologias que você não consegue defender em entrevista
%   - O multicols aqui é só visual --- o ATS lê linha por linha de
%     qualquer forma, então a ordem das categorias não importa para ele
%
% CATEGORIAS SUGERIDAS (adapte à sua área):
%   Linguagens | Backend/Framework | Arquitetura | Mensageria |
%   Banco de Dados | Testes | DevOps | Cloud | Versionamento | Métodos
%
% DICA ATS: Se a vaga menciona "Kafka" e você tem só RabbitMQ,
% não coloque Kafka. Mas se você tem ambos, liste os dois.
% Honestidade aqui evita problemas na entrevista técnica.
% =================================================================

\section*{Competências Técnicas}

\begin{multicols}{2}
\begin{itemize}[leftmargin=*, noitemsep]
  \item \textbf{Linguagens:} [Liste suas linguagens, ex: Java (8+), Python]
  \item \textbf{Backend:} [Frameworks e ferramentas, ex: Spring Boot, Spring MVC]
  \item \textbf{Arquitetura:} [Padrões, ex: Microsserviços, SAGA, Clean Code, SOLID]
  \item \textbf{Mensageria:} [ex: Kafka, RabbitMQ]
  \item \textbf{Banco de Dados:} [ex: PostgreSQL, MySQL (SQL); DynamoDB (NoSQL)]
  \item \textbf{Persistência:} [ex: JPA (Hibernate), Spring Data JPA]
  \item \textbf{Testes:} [ex: JUnit, Mockito]
  \item \textbf{DevOps:} [ex: Docker, Kubernetes, CI/CD (GitHub Actions)]
  \item \textbf{Cloud:} [ex: AWS (EC2, S3, Lambda, RDS, DynamoDB)]
  \item \textbf{Versionamento:} Git
  \item \textbf{Metodologias:} [ex: Scrum, Kanban]
\end{itemize}
\end{multicols}


% =================================================================
% BLOCO 4: EXPERIÊNCIA PROFISSIONAL
% =================================================================
% OBJETIVO: Provar que você já fez na prática o que a vaga pede.
% O ATS lê as palavras-chave; o recrutador lê os resultados.
% Você precisa satisfazer os dois.
%
% REGRAS GERAIS:
%   - Ordem cronológica inversa (mais recente primeiro)
%   - Cada item começa com VERBO DE AÇÃO no passado ou presente:
%     Desenvolveu / Implementou / Reduziu / Aumentou / Migrou /
%     Projetou / Liderou / Otimizou
%   - Nunca escreva "Responsável por..." --- isso descreve cargo,
%     não entrega. Descreva o que VOCÊ FEZ e qual foi o RESULTADO.
%   - Quantifique sempre que possível:
%     "Reduziu o tempo de resposta em 40%"
%     "Sistema processando 10k requisições/dia"
%     "Eliminou 100% dos processos manuais de X"
%   - Use as palavras-chave da vaga naturalmente nas frases
%
% ESTRUTURA DE CADA BULLET (fórmula CAR):
%   [Contexto/Ação] + [Tecnologia usada] + [Resultado/Impacto]
%   Ex: "Implementou comunicação assíncrona via Kafka entre 3
%   microsserviços, reduzindo o acoplamento e aumentando a
%   resiliência do sistema em cenários de falha parcial."
%
% SOBRE FREELANCE / AUTÔNOMO:
%   Nomeie os projetos reais sempre que possível.
%   "Sistema ERP para Empresa X (em produção desde 2023)" é
%   muito mais forte do que "Projeto de gestão para cliente".
%
% DICA ATS: Coloque o nome das tecnologias dentro das frases,
% não apenas na seção de competências. ATS pondera contexto.
% "Implementei APIs REST com Spring Boot" > só listar "Spring Boot"
% =================================================================

\section*{Experiência Profissional}

% -----------------------------------------------------------------
% COMO PREENCHER O CABEÇALHO DE CADA EXPERIÊNCIA:
%   \textbf{Cargo | Nome da Empresa ou Modalidade} \hfill Data início -- Data fim (ou Presente)
%
% Se for CLT: coloque empresa real
% Se for freelance: "Desenvolvedor Back-End | Freelance / Projetos para Clientes"
% Se for autônomo com projetos nomeados: mencione o projeto no bullet
% -----------------------------------------------------------------

\textbf{Cargo | Empresa ou Modalidade} \hfill mês/ano -- mês/ano (ou Presente) \\

\begin{itemize}[leftmargin=*, noitemsep]
  \item [Verbo de ação] + [o que foi feito] + [tecnologia usada] + [resultado ou escala].
  \item [Verbo de ação] + [padrão arquitetural aplicado] + [problema resolvido].
  \item [Verbo de ação] + [ferramenta/método] + [impacto mensurável se possível].
  % Repita para cada contribuição relevante. Ideal: 4 a 8 bullets por cargo.
  % Priorize o que a vaga pede. Se a vaga quer Kafka, seu bullet de
  % mensageria deve ser o primeiro ou segundo da lista.
\end{itemize}

\vspace{4pt}

% Repita o bloco acima para cada experiência anterior
% Experiências antigas e menos relevantes: 2 a 3 bullets apenas


% =================================================================
% BLOCO 5: PROJETOS RELEVANTES
% =================================================================
% OBJETIVO: Compensar pouco tempo de CLT com projetos reais.
% Para desenvolvedores com perfil freelance ou júnior/pleno,
% este bloco pode ser tão importante quanto a experiência.
%
% REGRAS:
%   - Liste apenas projetos com código no GitHub ou em produção
%   - Priorize: (1) sistemas em produção real, (2) projetos com
%     a stack da vaga, (3) projetos com resultado mensurável
%   - Cada projeto: nome + link + 2 a 3 bullets
%   - O primeiro bullet descreve O QUE É o sistema
%   - O segundo bullet descreve A STACK e decisões técnicas
%   - O terceiro bullet (opcional) descreve RESULTADO ou ESCALA
%
% DICA ATS: O nome do projeto não importa para o ATS.
% O que importa são as palavras-chave dentro dos bullets.
% Escreva "Backend Java com Spring Boot, APIs REST e JPA"
% em vez de apenas "Sistema de gestão".
%
% DICA RECRUTADOR: Projetos em produção real com nome de empresa/cliente
% têm muito mais peso do que projetos de portfólio genéricos.
% Se o sistema está rodando e sendo usado, diga isso explicitamente.
% =================================================================

\section*{Projetos Relevantes}

\textbf{Nome do Projeto --- Contexto (Ex: Em Produção / TCC / Open Source)} \hfill
\href{https://github.com/seu-usuario/repositorio}{github.com/seu-usuario/repositorio}

\begin{itemize}[leftmargin=*, noitemsep]
  \item [O que é o sistema e qual problema resolve --- contexto de negócio]
  \item [Stack técnica e decisões arquiteturais relevantes para a vaga]
  \item [Resultado, escala ou diferencial --- ex: "em uso contínuo desde X", "X usuários"]
\end{itemize}

\vspace{3pt}

% Repita para cada projeto relevante. Ideal: 2 a 4 projetos.
% Mais do que 4 começa a poluir o currículo.


% =================================================================
% BLOCO 6: FORMAÇÃO ACADÊMICA
% =================================================================
% OBJETIVO: Validação formal. ATS filtra por grau mínimo em algumas
% empresas. Coloque sempre o grau completo por extenso.
%
% REGRAS:
%   - Ordem: mais recente primeiro
%   - Use o nome oficial do curso e da instituição
%   - Para curso em andamento: deixe explícito "(em andamento)"
%     ou coloque a previsão de conclusão
%   - Não liste cursos técnicos ou de ensino médio aqui
%   - TCC relevante pode ser mencionado entre parênteses
%
% DICA ATS: Escreva o grau por extenso.
% "Bacharelado" e não apenas "Bach." ---
% alguns ATS não reconhecem abreviações.
% =================================================================

\section*{Formação Acadêmica}

\textbf{Grau (ex: Mestrado / Bacharelado) em Nome do Curso} -- Nome da Instituição (ano ou "em andamento")

\textbf{Grau em Nome do Curso} -- Nome da Instituição (ano de conclusão)


% =================================================================
% BLOCO 7: CURSOS E CERTIFICAÇÕES
% =================================================================
% OBJETIVO: Demonstrar atualização contínua e validar skills
% que não aparecem na formação formal.
%
% REGRAS:
%   - Liste certificações com nome oficial + instituição emissora
%   - Para certificações em andamento: coloque "(em andamento)"
%   - Certificações de cloud (AWS, GCP, Azure) têm alto peso
%     em vagas de backend moderno --- coloque em primeiro
%   - Não liste cursos sem certificado ou muito antigos (+5 anos)
%   - Não liste cursos irrelevantes para a vaga
%
% DICA ATS: Use o nome oficial completo.
% "AWS Certified Cloud Practitioner" tem mais peso nos filtros
% do que apenas "Cloud Practitioner" ou "AWS básico".
% =================================================================

\section*{Cursos e Certificações}

\textbf{Nome Oficial da Certificação} -- Instituição Emissora (ano ou "em andamento")

\textbf{Nome do Curso} -- Plataforma (ex: Udemy, Alura, Coursera)


% =================================================================
% CHECKLIST FINAL ANTES DE ENVIAR
% =================================================================
%
% [ ] O título do currículo espelha o título da vaga?
% [ ] As palavras-chave da vaga aparecem no resumo E na experiência?
% [ ] Cada bullet de experiência começa com verbo de ação?
% [ ] Pelo menos 2 bullets têm resultado mensurável (número, %,
%     escala, impacto)?
% [ ] Os projetos em produção estão nomeados com empresa/cliente?
% [ ] Não há tabelas, caixas de texto, colunas no meio do texto
%     principal ou ícones? (ATS não lê esses elementos)
% [ ] O arquivo será enviado como PDF gerado a partir do LaTeX?
%     (Nunca envie o .tex diretamente)
% [ ] O PDF tem menos de 2 páginas para vagas júnior/pleno?
%
% =================================================================

\end{document}